\documentclass[12pt]{article}
\usepackage[utf8]{inputenc}
\usepackage[spanish]{babel}
\usepackage{amsmath, amssymb, amsfonts}
\usepackage{geometry}
\geometry{a4paper, margin=2.5cm}
\usepackage{enumitem}

\title{Práctica en Clase: Problemas de Valor Inicial (PVI)}
\author{Ecuaciones Diferenciales}
\date{}

\begin{document}

\maketitle

\noindent\textbf{Instrucciones:} Resuelva los siguientes ejercicios detallando claramente su procedimiento. Dispone de 30 minutos para completar la actividad.

\vspace{0.8cm}

\begin{enumerate}[label=\textbf{\arabic*.}]
    \item % Ejercicio 1: Primer orden simple
    La función $y = c e^{-x} + 2x - 2$ es una familia uniparamétrica de soluciones de la ecuación diferencial $y' + y = 2x$. Encuentre la solución exacta del problema de valor inicial (PVI) de primer orden consistente en esta ecuación diferencial y la condición inicial $y(0) = 3$.

    \vspace{4.5cm}

    \item % Ejercicio 2: Segundo orden
    La función $y = c_1 e^x + c_2 e^{-x}$ es una familia biparamétrica de soluciones de la ecuación diferencial $y'' - y = 0$. Encuentre una solución del problema de valor inicial de segundo orden consistente en esta ecuación diferencial y las condiciones iniciales:
    \[ y(0) = 1, \quad y'(0) = 2 \]

    \vspace{6cm}

    \item % Ejercicio 3: Encontrar intervalo
    Dado que $y = \frac{1}{x^2 + c}$ es una familia uniparamétrica de soluciones de la ecuación diferencial de primer orden $y' + 2xy^2 = 0$.
    \begin{enumerate}[label=\textbf{\alph*)}]
        \item Determine una solución del problema de valor inicial que satisface la condición $y(2) = 1/3$.
        \item Indique el intervalo de definición $I$ más grande en el cual está definida la solución del inciso anterior.
    \end{enumerate}

    \vspace{6cm}

    \item % Ejercicio 4: Existencia y unicidad
    Considere el problema de valor inicial: 
    \[ \frac{dy}{dx} = \sqrt{y}, \quad y(x_0) = y_0 \]
    Utilice el Teorema de Existencia y Unicidad (Teorema de Picard) para determinar en qué regiones del plano $xy$ se garantiza que el problema tiene una solución única. ¿Qué ocurre específicamente si la condición inicial se da en el punto $(0,0)$?

\end{enumerate}

\newpage
\section*{Soluciones a los Ejercicios (Para el Profesor)}

\begin{enumerate}[label=\textbf{\arabic*.}]
    \item Sustituyendo $x=0$, $y=3$ en la familia de soluciones:
    \[ 3 = c e^{0} + 2(0) - 2 \implies 3 = c - 2 \implies c = 5 \]
    La solución del PVI es: $\mathbf{y = 5e^{-x} + 2x - 2}$.

    \vspace{0.5cm}

    \item Calculamos la derivada de la familia de soluciones:
    \[ y'(x) = c_1 e^x - c_2 e^{-x} \]
    Aplicamos las condiciones iniciales $y(0)=1$ y $y'(0)=2$:
    \begin{align*}
        1 &= c_1 + c_2 \\
        2 &= c_1 - c_2
    \end{align*}
    Sumando ambas ecuaciones: $3 = 2c_1 \implies c_1 = 3/2$. \\
    Restando la segunda ecuación a la primera: $-1 = 2c_2 \implies c_2 = -1/2$. \\
    La solución del PVI es: $\mathbf{y = \frac{3}{2}e^x - \frac{1}{2}e^{-x}}$.

    \vspace{0.5cm}

    \item
    \begin{enumerate}[label=\textbf{\alph*)}]
        \item Sustituimos $x=2$, $y=1/3$ en $y = \frac{1}{x^2 + c}$:
        \[ \frac{1}{3} = \frac{1}{2^2 + c} \implies 4 + c = 3 \implies c = -1 \]
        La solución es: $\mathbf{y = \frac{1}{x^2 - 1}}$.
        
        \item La función no está definida cuando el denominador es cero: $x^2 - 1 = 0 \implies x = \pm 1$.
        Los intervalos posibles donde la función es continua son $(-\infty, -1)$, $(-1, 1)$ y $(1, \infty)$.
        Como la condición inicial es en $x_0 = 2$, el intervalo más grande que contiene a este punto es $\mathbf{I = (1, \infty)}$.
    \end{enumerate}

    \vspace{0.5cm}

    \item Sea $f(x,y) = \sqrt{y}$. La función $f$ es continua para todo $y \ge 0$. \\
    Calculamos la derivada parcial respecto a $y$:
    \[ \frac{\partial f}{\partial y} = \frac{1}{2\sqrt{y}} \]
    Esta derivada parcial es continua para todo $y > 0$. \\
    Por el Teorema de Existencia y Unicidad, se garantiza una solución única en cualquier región rectangular del plano donde $\mathbf{y > 0}$ (es decir, en el semiplano superior, excluyendo el eje $x$). \\
    En el punto $\mathbf{(0,0)}$, la derivada parcial $\frac{\partial f}{\partial y}$ no es continua (ni siquiera está definida), por lo que el teorema \textbf{no garantiza unicidad}. (El problema admite múltiples soluciones por este punto, por ejemplo $y=0$ y $y=x^2/4$ para $x \ge 0$).
\end{enumerate}

\end{document}
